% Prediction chapter skeleton. Requires \TemplatePrompt from the selected paper family.
\section{预测模型的建立与检验}

\subsection{预测目标与信息边界}
\TemplatePrompt{定义预测对象、时间跨度、可用特征、信息截断点、评价指标和单位，防止未来信息泄漏。}

\subsection{可运行基线}
\TemplatePrompt{选择持久性、季节性朴素、历史均值或简单回归等同任务基线，并使用相同训练与测试切分。}

\subsection{模型结构与参数来源}
\TemplatePrompt{说明特征、损失函数、超参数、训练方案和模型选择理由；避免只罗列算法。}

\subsection{训练、滚动验证与复现}
\TemplatePrompt{记录时间切分、交叉验证方式、随机种子、软件版本、停止条件和产物哈希。}

\subsection{预测结果、区间与校准}
\TemplatePrompt{同时报告点预测、预测区间、基线、误差指标和单位；检查残差结构与区间覆盖率。}

\subsection{稳健性与结论边界}
\TemplatePrompt{检验窗口、特征、异常值和分布漂移的影响；说明模型在哪些时间或场景下不能外推。}
